%\begin{lemma}
%\end{lemma}
\so{\begin{lemma} \label{subexp bth}Let $\bth=\X^\dagger \y$ where $\y=\X\bts+\sigma \z$ where $\X,\z\distas\Nn(0,1)$. As $p,n\rightarrow \infty$, wpa.~1, there exists a constant $C>0$ such that, we have that
\[
\frac{1}{p}\sum_{i=1}^p|\sqrt{p}\bth_i|^k\leq C.
\]
\end{lemma}}
\begin{proof} We first show that entries of $\sqrt{p}\bth$ have finite $k$th moments for any $k\geq 0$. We split the proof in two components by writing $\bth=\ab+\sigma \bb$. 

First, let us focus on $\bb=\X^\dagger\z$. Write singular value decomposition $\X=\Ub\La\V^T$. Then $\bb=\V^T\La\Ub(\Ub\La^2\Ub^T)^{-1}\z=\V^T\La^{-1}\z'$ where $\z'=\Ub^T\z\distas\Nn(0,1)$. First, wpa.~1, $C\sqrt{p} \|\X\|\geq \sigma_{\min}(\X)\gtrsim c\sqrt{p}$ for some constants $C,c>0$ as the covariance $\bSi$ are bounded above and below by constants. To show subgaussianity, we fix unit length $\vb$ and consider the variable
\[
X=\vb^T\bb=\vb^T\V^T\La^{-1}\z'.
\]
by fixing all but $\z'$. Since $\z'$ is Gaussian, the subgaussian norm obeys 
\begin{align}
\tsub{\bb}=\sup_{\tn{\vb}=1}\tsub{X}&\lesssim \sup_{\tn{\vb}=1}\|\vb^T\V^T\La^{-1}\|_{\ell_\infty}\\
&\leq\sup_{\tn{\vb}=1} \tn{\vb^T\V^T}\sigma_{\min}(\X)^{-1}\\
&\leq \sigma_{\min}(\X)^{-1}\\
&\lesssim 1/\sqrt{p}.
\end{align}

To proceed, we address the subgaussianity of the main term 
\begin{align}
\ab&=\X^\dagger\X\bts=\X^T(\X\X^T)^{-1}\X\bts\\
&=\sqrt{\bSi}\Xb^T(\Xb\bSi\Xb^T)^{-1}\Xb\sqrt{\bSi}\bts\\
&=\sqrt{\bSi}\Qb^T (\Qb\bSi\Qb^T)^{-1}\Qb\sqrt{\bSi}\bts\\
&=\Vb^T\Vb\bts\\
&=\sum_{i=1}^n \vb_i\vb_i^T\bts.
\end{align} 
Here, $\Vb$ is the unitary matrix that is the range space of $\X$ and $\Qb$ is the unitary matrix that is the row space of $\Xb$. $(\vb_i)_{i=1}^n$ are rows of $\Vb$. Observe have that the range space of $\Vb$ is statistically identical to the range space of $\Qb\sqrt{\bSi}$. We will work with the properties of the row space $\Vb$ rather than the original design matrix $\X$. 

\noindent\textbf{Expectation of $\ab$:} We claim that $\E[\ab]=\m\odot \bts$ where $1\leq \m_j\geq 0$ for all $1\leq j\leq p$. Thus, by assumption on $\bts$, this would mean that $\E[|\sqrt{p}\E[\ab_i]|^k]$ and $p^{-1}\E[\|\sqrt{p}\E[\ab]\|_{\ell_k}^k]$ is bounded by a constant wpa.~1. To see this, observe that
\[
\E[\ab]=\sum_{i=1}^n\E[\vb_i\vb_i^T]\bts=\sum_{j=1}^p\sum_{i=1}^n\E[\vb_i\vb_i^T]\bts_j \eb_j,
\]
where $\eb_j$ is the $j$'th element of the standard basis. Note that the entries of $\vb_i$ have i.i.d.~Bernoulli signs same as the entries of the rows of $\Qb$. Thus $\E[\vb_i\vb_i^T]$ is a diagonal matrix $\Db^i$. This implies that, $\E[\X^\dagger\X]$ is diagonal and the $j$th index of $\bts$ is the only entry responsible from $\ab_j$. Thus
\[
\m_j=\sum_{i=1}^n \Db^i_{j,j}.
\]
Finally, $\sum_{i=1}^n \Db^i_{j,j}$ is equal to $\tn{\vb^\text{col}_j}^2\leq 1$ where $\vb^\text{col}_j$ is the $j$th column of $\Vb$ proving the result.

\noindent\textbf{Subexponentiality of $\ab$:} To conclude the proof, we show that $\ab$ is subexponential. Fix a unit length vector $\eb$ from standard basis. The claim is same as proving $\te{\sqrt{p}(X-\E[X])}\lesssim 1$ where
\begin{align}
X&=\eb^T\X^\dagger\X\bts=\eb^T\sqrt{\bSi}\Qb^T (\Qb\bSi\Qb^T)^{-1}\Qb\sqrt{\bSi}\bts.
\end{align}
\so{We will focus on the scenario $\bSi=\Iden_p$. General $\bSi$ is handled by the fact that $\bSi$ is upper and lower bounded by identity (thus isometric).} Thus, consider
\begin{align}
X&=\eb^T\X^\dagger\X\bts=\eb^T\Qb^T\Qb\bts\\
&=\frac{\tn{\Qb(\bts+\eb)}^2-\tn{\Qb(\bts-\eb)}^2}{4}.
\end{align}
Thus, our claim is same as showing
\[
\te{\tn{\Qb\vb}^2-\E[\tn{\Qb\vb}^2]}\lesssim 1/\sqrt{p}.
\]
for fixed $\vb=\bts+\vb$. Here, we can swap the randomness and assume fixed $\Qb$ and uniformly random $\vb$ entries of $\vb$ are equal to $\vb_i=\tn{\vb}\frac{\g_i}{\tn{\g}}$ for $\g\sim\Nn(0,\Iden_p)$ and $\tn{\vb}\lesssim 1$. Specifically, fixing $\Qb$ to be first $n$ elements of standard basis, we obtain
\[
\frac{1}{\tn{\vb}^2}(\tn{\Qb\vb}^2-\E[\tn{\Qb\vb}^2])=\frac{\sum_{i=1}^n \g_i^2}{\tn{\g}^2}-\frac{n}{p}=\frac{\sum_{i=1}^n \g_i^2-\frac{n}{p}\tn{\g}^2}{\tn{\g}^2}.
\]
To proceed, observe that both numerator and denominator are sub-exponential variables and concentrate accordingly to find that
\[
\Pro(|\tn{\Qb\vb}^2-\E[\tn{\Qb\vb}^2]|\gtrsim t/\sqrt{p})\leq \e^{-t}.
\]
Thus, entries of $\ab$ obeys $\te{X\sqrt{p}}\lesssim 1$ and together with bounded expectations, we obtain $\E[|\sqrt{p}X|^k]\lesssim 1$.

Finally, combining this with the $\bb$ term, we find the desired result $\E[|\sqrt{p}\bth_i|^k]\lesssim 1$.
%The final line follows from the fact that, since $\bSi_{i,i}$ are bounded below, wpa.~1, $\sigma_{\min}(\X)^{-1}\lesssim 1/\sqrt{p}$.

\so{To conclude with the proof, we use two things. First, $\frac{1}{p}\sum_{i=1}^p\bth_i$ weakly converges to the asymptotic limit given by Def.~\ref{def:Xi}. The proof is omitted for brevity but follows the exact same line of argument as the proof of Theorem \ref{thm:master_W2}. The main difference is that, for weak convergence proof, we don't require $k$th moments or pseudo-Lipschitz conditions (thus, the proof does not require this lemma and there is no circular logic). Once the weak convergence is established, we use Lemma 5 of \cite{bayati2011dynamics}. Since $k$th moments of $\bth_i$ are bounded and $\bth$ is converging weakly, this lemma implies that $\|\bth\|_{\ell_k}^k$ is bounded wpa.~1.}
\end{proof}



\cmt{
\[
\te{\tn{\Vb\eb}^2-\E[\tn{\Vb\eb}^2]}\lesssim 1/\sqrt{p}.
\]
}

\cmt{
Set $\A= (\Qb\bSi\Qb^T)^{-1/2}$. Observe that $\A\lesssim \Iden_n$ and that $\Vb=\A\Qb\sqrt{\bSi}$.
Set $\bar{X}=\eb^T\Qb^T\Qb\bts$. Observing $\Qb\bSi\Qb^T\gtrsim \Sigma_{\min}\Iden_{n}$, this yields that
\[
\E|X|^k\lesssim 
\]
Note that $\E[X]\lesssim 1/\sqrt{p}$ from above. Secondly, }









\cmt{
Let $\Qb$ be a uniformly random subspace satisfying $\Qb\Qb^T=\Iden_n$. Writing $\X=\bar{\X}\sqrt{\bSi}$ and noticing row space of $\bar{\X}$ is statistically identical to $\Qb$,
Now, let
%
\[
\A=\E[\Qb\bSi\Qb^T]^{-1/2}.
\]
Observe that, upper and lower eigenvalues of $\A$ are bounded by constants. Additionally, it follows from the properties of random tight frames that $\A$ is a diagonal matrix (specifically, using rows of $\Qb$ are symmetrically distributed). Now, observe that
\[
\A\Qb\bSi\Qb^T\A=\Iden.
\] 
Consequently $\A\Qb\sqrt{\bSi}\sim \Vb$. Thus, we can write
\[
\Vb^T\Vb\bt=\sqrt{\bSi}\Qb\A^2\Qb^T\sqrt{\bSi}\bt.
\]
We shall argue that, this vector is subexponential. First, given fixed vector $\vb$, we have that
\[
\E[\Qb\A^2\Qb^T\vb]
\]
we bound the infinity norm of the mean given by%$\|\E[\Vb^T\Vb\bt]\|_{\ell_{\infty}}$. Clearly
\[
\|\E[\Vb^T\Vb\bt]\|_{\ell_{\infty}}\leq 
\] }